\documentclass[10pt,a4paper]{amsart}
\usepackage[margin=19mm]{geometry}
\usepackage{amsmath,amssymb,mathtools}
\usepackage[scr=rsfs,cal=euler]{mathalfa}
\usepackage{xfrac}
\usepackage[unicode,hidelinks,bookmarks=false]{hyperref}
\newcommand{\ie}{i.e.\ }
\newcommand{\cf}{cf.\ }
\newcommand{\resp}{respectively\ }
\newcommand{\Subsection}{\subsection}
\providecommand{\nobreakdash}{\nobreak\hbox{-}}
\setlength{\emergencystretch}{2em}
\multlinegap0pt
\usepackage{bm}
\usepackage{bbm}
\usepackage{dsfont}
\usepackage{xspace}
\usepackage{cancel}
\usepackage{mathtools}
\usepackage{amsthm}
\makeatletter
\@ifundefined{theorem}{\newtheorem{theorem}{Theorem}}{}
\@ifundefined{lemma}{\newtheorem{lemma}[theorem]{Lemma}}{}
\makeatother
\title[A condition equivalent to the H\"{o}lder continuity of harmo]{A condition equivalent to the H\"{o}lder continuity of harmonic functions on unbounded Lipschitz domains}
\author{Marijan Markovic}
\date{Source: arXiv:2507.14511v1 (CC BY 4.0)}
\begin{document}
\maketitle

\noindent\emph{Source and licence.} A condition equivalent to the H\"{o}lder continuity of harmonic functions on unbounded Lipschitz domains. Version arXiv:2507.14511v1 (CC BY 4.0), distributed under the Creative Commons Attribution 4.0 International licence, as recorded in the arXiv licence metadata for this exact version and shown on the abstract page https://arxiv.org/abs/2507.14511v1. Full rights notice: licenses/arxiv-2507.14511-CC-BY-4.0.txt.\par\smallskip

\noindent\textit{Editorial scope.} Counted unit: the Lemma whose source label is \texttt{LE.DISTANCE.FROM.GRAPH} in the article (source0.tex lines 203--214, source label \texttt{LE.DISTANCE.FROM.GRAPH}), delivered together with the article's own complete original proof of it (source0.tex lines 216--249, one \texttt{proof} environment of 34 lines). No mathematics is written, repaired or re-derived: statement and proof are the article's own text line for line. Required context retained uncounted: none. Every definition, display and label of those lines is retained unchanged. Omitted from this package and not redistributed: 1--202, 215--215, 250--614. Nothing inside a retained excerpt is omitted or altered: every retained body is delivered exactly as the article's lines are written, so no omission receipt of any kind accompanies this package. Citations: the retained excerpts carry no citation command at all, so no bibliography entry of the article is transcribed and no reference is redistributed. Reference adaptation: none was needed and none was made; every reference inside the delivered unit resolves inside it, so no unresolved reference or citation is produced. Printed slips kept literal: no printed word, symbol, hypothesis or number of the counted statement or of its proof was altered. Layout and class: the article's own class is \texttt{amsart} with packages including amsmath, amsfonts, amssymb, amsthm, graphics; the wrapper is the lane's pinned \texttt{amsart} editorial wrapper at 10pt on A4, which restates the theorem-like environments the retained text needs and adds the article's own macro definitions that the retained text needs (none), with extra wrapper packages (bm, bbm, dsfont, xspace, cancel, mathtools). The delivered numbering is the unit's own. Assets: no retained span contains a figure environment, a graphics-inclusion command, a PDF-page inclusion, a table or a TikZ picture; no image file is copied into this package, the package declares no asset dependency and no image file is redistributed with it. Licence: version arXiv:2507.14511v1 of this article is distributed under the Creative Commons Attribution 4.0 International licence, as recorded in the arXiv licence metadata for this exact version and shown on the abstract page https://arxiv.org/abs/2507.14511v1. A full rights notice is delivered at licenses/arxiv-2507.14511-CC-BY-4.0.txt. Characters: every retained excerpt is pure ASCII, so the delivered engine, XeLaTeX with its default Latin Modern text font, reports no missing character.
\begin{lemma}\label{LE.DISTANCE.FROM.GRAPH}
Let $\Psi: \mathbb{R}^{N-1}\to\mathbb{R}$ be a $L$-Lipschitz function.   For $x = (x',x_N)\in E_\Psi$ and $\lambda\ge 0$
let us denote
\begin{equation*}
x_\lambda  =  x  + \lambda \mathbf {e}_N= (x', x_N + \lambda)\in E_\Psi.
\end{equation*}
Then we have
\begin{equation*}
d(x_\lambda, \partial E_\Psi) \ge d(x, \partial E_\Psi) \cos \gamma +  {\lambda\cos \gamma},
\end{equation*}
where  $\gamma  = \arctan L$.
\end{lemma}

\begin{proof}
Denote  $d  =  d(x, \partial E_\Psi)$   and   $d'= d(x_\lambda, \partial E_\Psi)$.   There exists $P\in \partial E_\Psi$
such that $ d = |x - P|$  ($P$ is any point where the distance is achieved).

Assume first that $P = (x', \Psi (x'))$. Let $C_\gamma$ be the cone with the vertex at  $P$,   with the axis parallel to
$\mathbf{e}_N$ and the opening $\frac {\pi}2-\gamma$.                          Denote $R = d(x, \partial C_\gamma) $ and
$R'= d(x_\lambda,  \partial C_\gamma)$.     Since $\Psi$ is Lipschitz,    we have $d'\ge  R'$.   From the similarity  of
triangles,  we obtain the following proportion
\begin{equation*}
R': R = (d+\lambda):d,
\end{equation*}
from which  we  find    $R'  = \frac R d (d+\lambda)$. Since $\frac Rd = \sin (\frac \pi 2 - \gamma) = \cos \gamma$,  it
follows
\begin{equation*}
d'\ge R'\ge (d +\lambda)\cos \gamma,
\end{equation*}
which we aimed to prove.

Let us now  consider the case where $P$ is not of the preceding form, that is, $P =(y', \Psi (y'))$, $y'\ne x'$. Then we
can    obtain a better estimate from below of $d(x_\lambda, \partial E_\Psi)$.  We consider the intersection of $E_\Psi$
with the plane  $\Pi$, which is determined by the vectors $\Vec{xP}$ and $\mathbf{e}_N$.  The intersection of  $\Pi$ and
the boundary of the cone $C_\gamma $ is made up of two straight lines that form the angle      $\frac \pi2-\gamma$  with
$\mathbf{e}_N$.  In  this case one of them, let it be denoted by $l_\gamma$,                is the tangent to the circle
$\partial B(x,d)\cap  \Pi$  at  $P$. Denote   $R' = d(x_\lambda, l_\gamma)$.  From the  similarity of  triangles,   this
time   we  obtain
\begin{equation*}
R': d  = \left(\frac d{\cos \gamma} +\lambda\right) : \frac d{\cos \gamma}.
\end{equation*}
From this proportion we have
\begin{equation*}
R' =  d + \lambda \cos \gamma\ge d \cos \gamma  + \lambda \cos\gamma,
\end{equation*}
which implies   the inequality in this  lemma, since   $d'\ge R'$.
\end{proof}

\end{document}
